% ============================================================================
% 09_conclusion.tex — Conclusion
% ============================================================================

\section{Conclusion}
\label{sec:conclusion}

This paper provides a first-time-buyer financial comparison of buying and renting-with-investing on English data, run over every monthly entry point from January 2005 to July 2021 with a five-year hold, on a representative dwelling whose price and rent share one regional composition. The answer to the popular question is that, over these two decades, there was close to no answer: buying produced greater terminal wealth in 102 of 199 cohorts (51\%), the median wealth ratio was 0.98, and the median gap \gbp{994} in the buyer's favour. Entry conditions (valuation and financing cost at the point of purchase) decided which side of parity a household landed on: renting won every cohort entering between 2006 and 2009 and again in 2016, buying won every cohort entering from 2011 to 2013 and in 2018 and 2019, and the regional buy-win share ran from 38\% in the North East to 63\% in the West Midlands.

That comparability extends well past five years. Buying won 62\% of seven-year holds and 61\% of ten-year holds, and the median renter still finished within 8--11\% of the owner, \gbp{14{,}454} behind at ten years on a dwelling worth \gbp{319{,}488}. Beyond five years buying is ahead more often than not, by margins small enough that renting-and-investing remains a defensible financial choice for a household that keeps to it. The distribution around that parity is asymmetric: buying won more cohorts, renting won the larger amounts. Renting more than doubled the owner's wealth in 21 of the 199 five-year cohorts and more than quadrupled it in five, all of which entered between August and December 2007, peaking at a ratio of 4.17 for the November 2007 cohort. That tail is why the mean ratio of 1.22 favours the renter while the mean gap of \gbp{1{,}744} favours the owner: renting's proportional wins landed where the owner's equity had collapsed to a small number, and buying's wins landed on much larger balance sheets.

At the short end the arithmetic is a judgement on transaction costs rather than on ownership. They consume roughly thirty per cent of a first-time buyer's initial outlay on the round trip and take about six years to amortise. Buying won only 30\% of one-year holds, so a household expecting to move within two or three years is, on this evidence, usually better off renting. The shape of that profile belongs to this sample, which contains one crash, one long expansion of cheap credit and one exceptional equity bull market, and I would not read it as a general fact about tenure.

Beyond the headline, the analysis yields a methodological caution and a forward-looking metric. The caution is that the official national house-price and rent aggregates weight regions differently, so their ratio (the most widely quoted price-to-rent figure) lies outside the range of every English region and flatters ownership in any cost comparison built on it. The metric is the required appreciation rate. At June 2026 conditions a five-year buyer needs 2.3\% annual appreciation to match a cautiously parameterised global equity portfolio and 3.2\% to match the return that portfolio actually delivered, while the dwelling itself appreciated at 3.2\%. The forward-looking arithmetic is therefore genuinely balanced, in the specific sense that plausible assumptions fall on both sides of it. Reporting only the cautious figure would make the case for buying look settled when it is not.

The results also discipline how the trade-off should be described. Leverage is what keeps buying competitive, rather than house-price growth: an unlevered cash buyer wins only 8\% of cohorts. But more leverage is not simply better. Buying does best at deposits of 15--20\% and worse at 5\%, because leverage magnifies the losses of a badly-timed purchase more than it magnifies the gains of a well-timed one. In the North East and North West, sixteen cohorts entering between mid-2007 and mid-2008 ended their five years with negative net worth after selling costs. Conversely, the renter's parity is earned by conduct: it requires investing the full cost difference in a \emph{globally} diversified portfolio, and a renter who bought only UK equities loses. Two calibrations move the answer more than any others: the maintenance-and-depreciation charge and the fixation design. Section~\ref{sec:robustness} reports the full mapping from each to the headline so that readers can relocate the result under their own priors.

All data, code, and an interactive implementation of the framework are openly available, and the analysis is designed to be re-run as new months of official data arrive. Extensions to Scotland and Wales as their rent statistics mature, first-time-buyer-specific price series matched to an appropriate rent, and the related buy-to-let decision are natural next steps.
